\subsection{Umrechnen der Daten}
Entsprechend \cref{table:FehlerVM} wurden die im Messprotokoll erhobenen Daten umgerechnet, mit signifikanten Stellen zu den Tupeln \(U\pm\Delta U\) verarbeitet. 
Mit der Fehlerformel 
\begin{equation}
    \Delta \sqrt{U-U_0}=\sqrt{\left(\frac{1}{2\sqrt{U-U_0}} \Delta U\right)^2+\left(\frac{-1}{2\sqrt{U-U_0}} \Delta U_0\right)^2}
\end{equation}
ergeben sich die Werte, die im Anhang angehängt sind.
Dort finden sich auch die entsprechenden Diagramme, die mit dem Programm LibreOfficeCalc erstellt worden sind. Die händisch erstellten Diagramme sind immer in entsprechenden Subsections eingefügt.

\subsection{Sperrspannung}
Die aufgenommenen Werte aus Tabelle \secref{sec:umgerechnetedaten} wurden zur Erstellung von linearen Fits genutzt, um anhand der Trendlinie, die über die linearen, mittleren Bereiche gelegt wurde, den Wert für die Sperrspannung \(U_S\) zu ermitteln. \\Man sieht im Diagramm \cref{sec:allefarben}, dass die Fehlerbalken minimal sind, dementsprechend konnte keine Fehlergerade erstellt werden. Aus diesem Grund wurde der Fehler jeder Sperrspannung aufgrund der Nichtlinearität der Werte auf \(\Delta U_S=\pm \qty{0.1}{\volt}\) nach oben abgeschätzt. Es ergibt sich folgende Tabelle:
\emph{Anmerkung:} Mittlerweile habe ich es auch händisch erstellt, als x-Fehler (für die Gegenspannung \(U)\) habe ich mit \(\Delta U=\qty{0.01}{\volt}\) gearbeitet. Für die nur schwer einzeichbaren, nach Augenmaß (also vermutlich beschränkte  Authentizität ) erstellten Fehlergeraden wurde auch der Schnitt mit der x-Achse bestimmt, \(\Delta U_S\) ergab sich dann aus der Differenz der zwei Schnittpunktwerte. Die Fehler für \(\sqrt{U_I-U_{I_0}}\) sind marginal (maximal \(\qty{0.004}{\volt}\equiv \qty{0.4}{\milli\meter}\leftarrow\) das ist die Größe auf den Zeichnungen des y-Fehlerbalkens), weswegen sie nicht eingezeichnet wurden. Die Werte finden sich jetzt auch in der folgenden Tabelle.

\begin{table}[htbp]
    \centering 
    \caption{Aus \cref{sec:allefarben} extrahierte Sperrspannungen, sowie Zeichnungen per Hand \cref{sec:haendisch}}
\begin{tabularx}{\textwidth}{cxZZZ}
\toprule
Farbe & \nu [\unit{\hertz}] & U_{S\;LBC} [\unit{\volt}] & U_{S\;H} [\unit{\volt}] &S\;[\sigma] \\ \midrule
UV       & 821.3  & \qty{-2.02(10)}{} & \qty{-2.06(5)}{} & 0.8 \\
Violett & 740.2  &  \qty{-1.59(10)}{} & \qty{-1.60(2)}{} & 0.5\\
blau     & 687.9  &  \qty{-1.53 (10)}{} & \qty{-1.47(8)}{} & 0.8\\
grün     & 549.0  &  \qty{0.98(10)}{} & \qty{-0.98(6)}{}& 0 \\
gelb     & 518.7  &  \qty{-0.76(10)}{} & \qty{-0.750(20)}{} & 0.5\\ \bottomrule
\end{tabularx}
\label{table:us}
\end{table}

Die Abweichung hält sich in Grenzen, die Sigma-Abweichung wurde ohne Einbeziehen vom abgeschätzten Fehler bestimmt, denn dieser würde das ganze nochmal ordentlich runterdrücken.
\begin{equation}
    \sigma_{Abw}=\frac{\mid U_{S,\;LBC}-U_{S\;H} \mid}{\sqrt{\Delta U_{S\;H}^2}}
\end{equation}

\clearpage

\markright{LBC: Bestimmung der Sperrspannung}
\includepdf[pages=-,landscape=true,pagecommand={}, addtotoc={1, subsubsection, 3, LBC: Bestimmung der Sperrspannung, sec:allefarben}]{allefarben.pdf}

\markright{Händisch: Bestimmung der Sperrspannung}
\includepdf[pages=1-5,pagecommand={}, addtotoc={1, subsubsection, 3, Händisch: Bestimmung der Sperrspannung, sec:haendisch}]{Diagramme.pdf}

\markright{\thesection. Auswertung}
\subsection{Plancksches Wirkungsquantum}
Aus \cref{table:us} wurde \cref{fig:sperrspannung.png} erstellt:
\xfigure{htbp}{scale=0.9}{sperrspannung.png}{Diagramm Sperrspannnug}{\(U_S(\nu)\) Ausgleichsgerade und Fehlergerade (schwarz)}

Nun kann nach \cref{eq:hs} aus der Steigung der \(U_S(\nu)\) Trendlinie \(m=\frac{h}{e}\) das Planksche Wirkungsquantum \(h\) berechnet werden.
Aus Steigung der Ausgleichs- und Fehlergeraden ergibt sich: 
\[m=\qty{0.0039(9)}{\volt\per\tera\hertz}\]
Der Fehler für \(h\) wird folgend berechnet:
\begin{equation}
    \Delta h=e\Delta m
\end{equation}
Final steht das Ergebnis:
\begin{rb}
h=\qty{6.2(15)e-34}{\joule\s}
\end{rb}

Als Ergebnis für das zeichnerische Ermitteln der Sperrspannung aus \cref{sec:Sperrspannung} ergibt sich:
\[m=\qty{0.0041(5)}{\volt\per\tera\hertz}\]
Damit berechnet sich:
\begin{rb}
h=\qty{6.6(8)e-34}{\joule\s}
\end{rb}

\clearpage
\markright{Zeichnerische Auftragung von Sperrspannung gegen Frequenz}
\includepdf[pages=6,pagecommand={}, addtotoc={6, subsubsection, 3, Zeichnerische Auftragung von Sperrspannung gegen Frequenz, sec:Sperrspannung}]{Diagramme.pdf}



% \begin{table}[htbp]
%   \centering
%   \caption{}
%   \begin{tabularx}{\textwidth}{ZZZZ}
%     \toprule
%             \cmidrule[0.3pt](l{1.0cm}r{1.5cm}){1-4}
%     \bottomrule
%   \end{tabularx}
% \label{table:}  
% \end{table}
